\subsection*{Part 4: Predict it!}
\begin{center}
\renewcommand{\arraystretch}{1.3}
\begin{tabular}{l l P{7.5cm}}
\toprule
\textbf{Code} & \textbf{Shows} & \textbf{Why} \\
\midrule
\py|print(9 / 3)| & \texttt{3.0} & Divide always gives a decimal. \\
\py|print("4" + "4")| & \texttt{44} & They are strings, so \texttt{+} joins them. \\
\py|print("Dog" == "dog")| & \texttt{False} & Capital letters count in Python. \\
\py|print("pizza"[1])| & \texttt{i} & Python counts from 0, so \texttt{[1]} is the second letter. \\
\py|print(round(4.5))| & \texttt{4} & A .5 goes to the nearest even number. \\
\py|print(17 % 5)| & \texttt{2} & $17 = 3 \times 5 + 2$, so 2 is left over. \\
\bottomrule
\end{tabular}
\end{center}

\subsection*{Part 4, Mission 1: Dice roller}
For 100 rolls, change line 3 to \py|for i in range(100):| and line 8 to \py|str(total / 100)|. With two dice, the biggest total is $6 + 6 = 12$.

\subsection*{Part 4, Mission 2: Times-table printer}
For the 9 times table, change line 1 to \py|table = 9|. To go up to 20, change line 3 to \py|range(20)|. For square numbers, line 4 becomes \py|print(f"{n} x {n} = {n * n}")|. The \texttt{True} or \texttt{False} at the end tells you whether the answer is even.

\subsection*{Part 4, Mission 3: Random dot art}
With 7 colours, line 15 needs \py|random.randint(0, 6)|. For the top half, line 12 becomes \py|y = random.randint(0, 180)|.

\subsection*{Part 4, Mission 4: Name machine}
It says \py|i + 1| because \py|i| starts at 0, but people count letters from 1.

\subsection*{Part 4: Translate it yourself}
Your code may look a little different and still be right. If it prints three totals, you did it!
\begin{codenonum}
import random
for i in range(3):
    a = random.randint(1, 10)
    b = random.randint(1, 10)
    print("Total: " + str(a + b))
\end{codenonum}
One real run printed \texttt{Total: 9}, \texttt{Total: 19} and \texttt{Total: 8}. An f-string works too: \py|print(f"Total: {a + b}")|.

\subsection*{Part 4: Two-line puzzle}
\begin{codenonum}
word = "Turtle"
print(word[0] + str(len(word)))
\end{codenonum}
\py|word[0]| is the first letter, \texttt{T}. \py|len(word)| is 6, and \py|str( )| turns it into a string so it can be joined.
